% ============================================================
% sec_toroidal_triad.tex
% W33-Theory — Section: The Toroidal Triad
% Szilassi (×2), Császár (×5), and the Tetrahedron
% ============================================================

\section{The Toroidal Triad: Szilassi, Cs\'asz\'ar, and the Tetrahedron}
\label{sec:toroidal_triad}

Three polyhedra form a single combinatorial family related by duality,
a common genus-0 seed, and the Fano self-duality of $\mathrm{PG}(2,2)$.
Their invariants are collected in Table~\ref{tab:triad}.

\begin{table}[h]
\centering
\begin{tabular}{lccccl}
\toprule
Polyhedron & $V$ & $E$ & $F$ & Genus & Role \\
\midrule
Tetrahedron & 4 & 6 & 4 & 0 & Seed: satisfies both hole equations at $h=0$ \\
Cs\'asz\'ar (\texttimes5) & 7 & 21 & 14 & 1 & $K_7$ triangulation of $T^2$; minimal triangulation \\
Szilassi (\texttimes2) & 14 & 21 & 7 & 1 & Heawood dual; 7-colour lower bound \\
\bottomrule
\end{tabular}
\caption{The toroidal triad. Edge count $E=21=\binom{7}{2}$ is invariant under the $V\leftrightarrow F$ duality.}
\label{tab:triad}
\end{table}

\subsection{The Dual Hole Equations}

Two Euler-characteristic equations, one for each adjacency type,
give the set of polyhedra in which every pair of vertices (resp.\ faces)
is connected by an edge (resp.\ share an edge):
\begin{equation}
  h_{\mathrm{vertex}}(v) = \frac{(v-3)(v-4)}{12}, \qquad
  h_{\mathrm{face}}(f)   = \frac{(f-4)(f-3)}{12}.
  \label{eq:hole_eqs}
\end{equation}
Both equations admit integer solutions at exactly the same residues
$n \equiv 0,3,4,7 \pmod{12}$.
The first two non-trivial solutions are $n=4$ ($h=0$, the tetrahedron)
and $n=7$ ($h=6$, the torus polyhedra).  The next predicted level is $n=12$
(resp.\ $n=44$), with invariant $E = \binom{12}{2} = 66$ edges.

\subsection{The Tetrahedron as Intermediary}

The tetrahedron is the \emph{only} known polyhedron satisfying
\emph{both} hole equations simultaneously ($h=0$, $v=f=4$).
It is also the only other polyhedron besides Szilassi in which every
face shares an edge with every other face, \emph{and} the only other
manifold besides Cs\'asz\'ar with no diagonals.
It therefore sits precisely between the two toroidal types, mediating
both vertex-adjacency and face-adjacency simultaneously.

\subsection{Fano Plane as Self-Dual Bridge}

The Fano plane $\mathrm{PG}(2,2)$ unifies both:
reading its 7 points as \emph{vertices} yields Cs\'asz\'ar;
reading its 7 points as \emph{faces} yields Szilassi.
Since $\mathrm{PG}(2,2)$ is self-dual, swapping points $\leftrightarrow$ lines
is isomorphic to the $V\leftrightarrow F$ duality between the two polyhedra.
The flag-orbit count $42 = 6\times 7$ shared by both polyhedra
equals $30 + 12$ (600-cell contribution plus the 12-periodic lattice)
and encodes $b_0^{\mathrm{QCD}} = \Phi_6 = 7$.

\subsection{Minimal Triangulations}

The Heawood bound gives the minimum number of vertices needed to
triangulate a surface of genus $g$:
\begin{equation}
  V_{\min}(g) =
  \begin{cases} 4 & g=0, \\
  \left\lceil\dfrac{7+\sqrt{1+48g}}{2}\right\rceil & g\geq 1.
  \end{cases}
  \label{eq:heawood}
\end{equation}
The tetrahedron achieves the bound for $g=0$ ($V=4$),
and Cs\'asz\'ar achieves it exactly for $g=1$ ($V=7$),
making it the \emph{unique} minimal triangulation of the torus.
The difference $\Delta V = 7-4 = 3 = \binom{3}{2}$ mirrors the
3-simplex step and the mod-12 residue shift from $4\to 7$.

\subsection{Realization Counts and $\Phi_6=7$}

The Szilassi polyhedron has exactly \textbf{2 realizations}
(related by its $\mathbb{Z}_2^+ \times I$ espyric symmetry)
and Cs\'asz\'ar has exactly \textbf{5 realizations}
(M\"obius 1858; Cs\'asz\'ar 1949; Bokowski--Eggert 1991 $\times 3$).
Their sum $2 + 5 = 7 = \Phi_6$ is the QCD beta-function coefficient
$b_0 = 11 - \tfrac{2}{3}n_f$ at $n_f = 6$ and simultaneously
the first nontrivial integer solution of both hole equations,
the order of the Fano plane, and the chromatic number of the torus.
